% !TeX root = ../../../main.tex


In this appendix, we provide detailed descriptions of data generation processes, methodological explanations, experimental setups, and results supplementing the main text. We begin by describing the mathematical formulations used to create our synthetic datasets, followed by detailed explanations of imputation methods and causal inference techniques. We then provide comprehensive information about model training procedures and hyperparameter settings for reproducibility. The appendix concludes with additional experimental results on synthetic, semi-synthetic, and real-world datasets.
% This appendix provides supplementary material supporting the main text, including detailed descriptions of data generation processes, methodological explanations, experimental setups, and additional results. We begin by describing the mathematical formulations used to create our synthetic datasets, followed by detailed explanations of imputation methods and causal inference techniques. We then provide comprehensive information about model training procedures and hyperparameter settings for reproducibility. The appendix concludes with extensive additional experimental results on synthetic, semi-synthetic, and real-world datasets.
%, as well as a discussion of broader impacts.

We would like to declare \textbf{the use of Large Language Models (LLMs)} in this work. LLMs were used as general-purpose assistive tools. Specifically, they supported parts of the writing process (editing, formatting, and polishing text) without contributing to the core methodology, scientific rigor, or originality of the research. In addition, LLMs were used to assist with improving visualization code for figures, documenting the code, and minor refactoring. No part of the conceptualization, design, or execution of the research relied on LLMs.

\textbf{Appendix~\ref{ch04:app:synthetic-dataset}: Additional Details of the Synthetic Datasets.}
This section provides the mathematical formulations for generating covariates, treatment assignments, event times, and censoring times across different scenarios. It describes how the synthetic datasets systematically vary across causal configurations and survival scenarios, including details on covariate generation, treatment assignment mechanisms, event time generation, censoring time generation, and observed data construction.
This section also includes Kaplan-Meier curves for the synthetic event-time and censoring distributions, illustrating scenario-level variation used throughout the benchmark.

\textbf{Appendix~\ref{ch04:app:imputation}: Imputation Methods Details.}
This section explains three surrogate imputation strategies for estimating true event time in right-censored survival data: Margin Imputation, IPCW-T Imputation, and Pseudo-observation Imputation. It provides mathematical formulations for each method and discusses their respective advantages and limitations.

\textbf{Appendix~\ref{ch04:app:exp-setup}: List of CATE Estimators in \textsc{SurvHTE} Benchmark.}
This section details the \ChFourRebuttal{53} different conditional average treatment effect (CATE) estimator variants evaluated in the benchmark, including outcome imputation methods, direct-survival CATE models, and survival meta-learners, with a breakdown of how these variants are constructed.

\textbf{Appendix~\ref{ch04:app:causal_method_details}: Detailed Overview of Causal Inference Methods.}
This section provides comprehensive explanations of various causal inference methods, including meta-learners (T-Learner, S-Learner, X-Learner, DR-Learner), Double-ML, Causal Forest, Causal Survival Forests, SurvITE, and Survival Meta-Learners, discussing their implementation in survival contexts.

\textbf{Appendix~\ref{ch04:app:method_settings}: Model Training Details and Hyperparameters.}
This section covers the hyperparameter grids, model selection procedures, and computational costs associated with each method class evaluated in the benchmark, providing details on the experimental setup for reproducibility.

\textbf{Appendix~\ref{ch04:app:additional_synthetic_exps}: Additional Experimental Results for Synthetic Dataset.}
This section presents comprehensive experimental results on synthetic datasets, including full rankings of models, win-rate summaries (Top-1 / Top-3 / Top-5) across methods, performance across different survival scenarios and causal configurations, detailed CATE RMSE and ATE bias plots, evaluation of auxiliary components, and convergence behavior under varying training set sizes.

\textbf{Appendix~\ref{ch04:app:more-semi-syn-data-setup-result}: Semi-Synthetic Datasets.}  
This section includes data setup and detailed analysis of semi-synthetic datasets derived from ACTG 175 and MIMIC-IV, including covariate statistics, censoring rate range, and comprehensive performance results across methods. It also presents results for the survival-probability-based CATE estimand across multiple horizons, and includes sensitivity analyses of the RMST-based CATE where the evaluation horizon is varied.

\textbf{Appendix~\ref{ch04:app:more-real-data-setup-result}: Real-World Datasets.}
This section provides detailed descriptions of data preprocessing and additional experimental results for the Twins dataset and the ACTG 175 HIV clinical trial dataset, including CATE RMSE results with different time horizons and comparisons of CATE estimates between baseline and high-censoring conditions.

\textbf{Appendix~\ref{ch04:app:informative-censoring-uconf}: Informative Censoring via Unobserved Confounding.}
An additional dataset violating the ignorable censoring assumption via latent confounders. This illustrates extensibility of \textsc{SurvHTE-Bench} to incorporate alternative censoring mechanisms beyond the $8 \times 5$ design.
% Informative Censoring via Unobserved Confounding.} An additional dataset illustrating how censoring can depend on latent variables, violating the ignorable censoring assumption. Demonstrates the extensibility of \textsc{SurvHTE-Bench} to incorporate alternative censoring mechanisms beyond the $8 \times 5$ design.


% \paragraph{Appendix~\ref{app:broader_impacts}: Broader Impacts.}
% This section discusses the potential positive impacts of the benchmark for improving personalized medicine and clinical decision-making, while also acknowledging potential risks if methods are misapplied in practice, emphasizing the importance of careful validation and human oversight in clinical settings.




